\begin{textsample}{NEG Dim 4 – 2023\_c5 – Score -1.00 – 2023\_c5/enade\_2023\_nutri\_ab5bde7002.txt}  \label{ex:f4_neg_097}
CAPÍTULO 4 DISTRIBUIÇÃO DE CONCEITOS A finalidade deste capítulo é apresentar um panorama nacional sobre a distribuição do Conceito Enade 2023.
Para traçar esse panorama, recorre -se a tabelas e gráficos, que são acompanhados de comentários e análises que contribuem para a contextualização dos dados neles expostos.
Para visualizar o panorama do Enade 2023, em nível nacional, na primeira seção deste capítulo, os dados sobre o número de cursos participantes são organizados por grande região.
Na segunda, o mapeamento do conceito do exame é feito por categoria administrativa das IES e modalidade de oferta dos cursos, de maneira cruzada às grandes regiões.
Na última seção, a lógica de apresentação dos dados é feita por organização acadêmica e, também, por grande região de funcionamento dos cursos, no país como um todo. 4.1 PANORAMA NACIONAL DA DISTRIBUIÇÃO DOS CONCEITOS A lógica adotada na apresentação dos dados, nesta seção e nas duas seguintes, pressupõe o registro de duas informações de natureza reiterativa.
O primeiro diz respeito aos valores possíveis do Enade, que variam em níveis de 1 a 5, sendo que, à medida que esses valores aumentam, o curso é melhor avaliado, e ao fato de que aqueles cursos que tiveram somente um ou nenhum concluinte participante não obtêm o Conceito Enade, ficando sem conceito ( SC ), como ressaltado na apresentação deste relatório.
O segundo registro é que, caso haja diferença entre os cursos tabulados neste capítulo e no capítulo 1, ela corresponderá aos cursos que não foram avaliados, ou seja, àqueles, em princípio, sem estudantes concluintes que tivessem sido inscritos para a prova.
Na Tabela 4.1, são apresentadas a quantidade e a distribuição de cursos de Nutrição participantes do Enade 2023, por faixa de conceito e grande região.
A partir dos dados da Tabela 4.1, nota -se que, dos 621 cursos participantes, 253 ( 40,7 \% ) classificaram -se com conceito 3, o valor modal.
Com exceção da região Norte ( 26,5 ), esse foi também o conceito modal em todas as regiões : Centro-Oeste ( 36,4 \% ), Nordeste ( 37,5 \% ), Sudeste ( 43,9 \% ), Sul ( 44,0 \% ).
O conceito 4 foi o segundo mais frequente em nível nacional ( 25,1 \%, correspondendo a 156 cursos ), tendo sido o segundo conceito mais frequente na maioria das regiões, exceto no Nordeste e no Norte.
O conceito 2 foi o terceiro ( 21,3 \%, correspondendo a 132 cursos ).
Houve, ainda, 48 ( 7,7 \% ) cursos que receberam o conceito 5 e 12 ( 1,9 \% ) cursos que receberam o conceito 1.
Dos 621 cursos de Nutrição, 20 ( 3,2 \% ) ficaram sem conceito ( SC ).
A região Norte participou com 34 cursos, o que corresponde a 5,5 \% do total nacional.
Desses, 16 cursos receberam o conceito 2, o que equivale a 47,1 \% do total regional, o conceito modal na região.
A nove cursos ( 26,5 \% ) atribuiu -se o conceito 3 ; a seis cursos ( 17,6 \% ) o conceito 4 ; e a dois cursos ( 5,9 \% ) o conceito 1.
Ainda nessa região, nenhum curso ( 0,0 \% ) obteve conceito 5 e um curso ( 2,9 \% ) ficou sem conceito ( SC ).
A região Nordeste participou com 152 cursos, correspondentes a 24,5 \% do total nacional.
Nessa região, 57 cursos ( 37,5 \% em termos regionais ) obtiveram conceito 3, o conceito modal para a região, como já comentado.
Os conceitos 2 e 4 foram atribuídos, respectivamente, a 47 e a 24 cursos ( 30,9 \% e 15,8 \% ).
Ainda nessa região, 13 cursos ( 8,6 \% ) foram avaliados com o conceito 5 e três cursos ( 2,0 \% ) foram avaliados com o conceito 1.
Outros oito cursos ( 5,3 \% ) ficaram sem conceito ( SC ).
Para os 280 cursos participantes da região Sudeste ( 45,1 \% do total nacional ), o conceito 3 foi o conceito modal, como já anteriormente assinalado, tendo sido obtido por 123 cursos ( 43,9 \% ).
O conceito 4 foi atribuído a 78 cursos ( 27,9 \% ) e o conceito 2, a 47 cursos ( 16,8 \% ).
Receberam os conceitos 5 e 1, respectivamente, 22 cursos ( 7,9 \% ) e quatro cursos ( 1,4 \% ).
Nessa região, seis cursos ( 2,1 \% ) ficaram sem conceito ( SC ).
Os 100 cursos da região Sul corresponderam a 16,1 \% do total nacional.
Como já indicado, houve predominância do conceito 3 ( conceito modal ), atribuído a 44,0 \% dos cursos da região, ou seja, atribuído a 44 dos 100 cursos participantes na região Sul.
O conceito 4 foi atribuído a 37 cursos ( 37,0 \% ) ; o conceito 5, nove cursos ( 9,0 \% ) ; e o conceito 2, a seis cursos ( 6,0 \% ).
Na região Sul, um curso ( 1,0 \% ) obteve conceito 1 e três cursos ( 3,0 \% ) ficaram sem conceito ( SC ).
Já dos 55 cursos participantes na região Centro-Oeste ( 8,9 \% do total nacional ), 20 ( 36,4 \% em termos regionais ) receberam conceito 3, o conceito modal, como já destacado. 16 cursos ( 29,1 \% ) obtiveram o conceito 2, 11 cursos ( 20,0 \% ) obtiveram o conceito 4 e quatro cursos ( 7,3 \% ) o conceito 5.
Nessa região, dois cursos ( 3,6 \% ) obtiveram o conceito 1 e dois cursos ( 3,6 \% ) ficaram sem conceito ( SC ).
Em um gráfico de distribuição cumulativa, em cada valor de abscissa ( eixo x ), o valor apresentado no eixo das ordenadas ( eixo y ) é o valor acumulado da variável ( Conceito Enade ).
Em linhas gerais, em um gráfico desta mesma tipologia, poligonais mais à esquerda representam, no caso dos dados apresentados nesta seção, as grandes regiões com a pior distribuição, e poligonais mais à direita, as grandes regiões com os melhores conceitos.
A partir destas diretrizes, no Gráfico 4.1, é apresentada a distribuição cumulativa do Conceito Enade, segundo a grande região.
No Gráfico 4.1, observando a linha vertical pontilhada, que ajuda a visualizar o percentual de conceitos acumulados até o valor 3, ou seja, conceitos 1, 2 e 3, pode -se notar, por exemplo, que a região Norte ( linha azul ) apresenta 81,8 \% dos cursos com conceito 3 ou abaixo deste valor.
A região Sul ( linha roxa ), por outro lado, apresenta 52,6 \% dos cursos na mesma situação : com conceito 3 ou abaixo.
Destaca -se que, na região Norte, o conceito 5 não foi alcançado por nenhum curso, na qual todos os cursos têm conceito menor ou igual a 4, e a respectiva poligonal cumulativa alcança 100 \% no conceito 4.
A região Norte apresenta os cursos com conceitos mais baixos, ao passo que a região Sul apresenta a distribuição com valores maiores.
Os cursos em IES situadas nas regiões Centro-Oeste e Nordeste ( linhas laranja e verde, respectivamente ) apresentam uma situação intermediária. 4.2 CONCEITOS POR CATEGORIA ADMINISTRATIVA, MODALIDADE DE OFERTA DOS CURSOS E GRANDE REGIÃO Na Tabela 4.2, é apresentada a distribuição dos cursos de Nutrição participantes do Enade 2023, por categoria administrativa e modalidade de oferta, de acordo com os conceitos por eles alcançados, segundo as grandes regiões brasileiras.
De acordo com as informações da Tabela 4.2, dos 621 cursos participantes, 72 ( 11,6 \% ) eram ministrados em instituições públicas e 549 ( 88,4 \% ) em instituições privadas.
Quanto à modalidade de oferta, 573 ( 92,3 \% ) cursos eram presenciais e 48 ( 7,7 \% ), a distância.
De acordo com as informações da mesma tabela, em termos nacionais, entre instituições públicas e privadas, observa -se que, dos 48 cursos avaliados com conceito 5, 33 foram oferecidos em IES públicas, e 15, em IES privadas.
Dos 549 cursos participantes de IES privadas, o conceito 3 foi o valor modal, atribuído a 245 ( 44,6 \% ) cursos.
Em relação aos demais cursos de instituições privadas, 129 cursos ( 23,5 \% ) foram avaliados com o conceito 2, 128 cursos ( 23,3 \% ) foram avaliados com o conceito 4, 15 cursos ( 2,7 \% ) foram avaliados com o conceito 5 e 12 cursos ( 2,2 \% ) com o conceito 1.
Nessa categoria, 20 cursos ( 3,6 \% ) ficaram sem conceito ( SC ).
Na rede pública, o conceito modal foi o 5, atribuído a 33 cursos, correspondentes a 45,8 \% dos 72 cursos da categoria.
Dos demais cursos participantes, 28 cursos ( 38,9 \% ) obtiveram o conceito 4, oito cursos ( 11,1 \% ) receberam o conceito 3 e três cursos ( 4,2 \% ) o conceito 2.
O conceito 1 foi atribuído a nenhum curso ( 0,0 \% ).
Nessa categoria administrativa, nenhum curso ( 0,0 \% ) ficou sem conceito ( SC ).
Dos 573 cursos participantes da modalidade presencial, o conceito 3 foi o valor modal, atribuído a 226 ( 39,4 \% ) cursos.
Dos demais cursos dessa modalidade de oferta, 147 cursos ( 25,7 \% ) receberam o conceito 4 ; 123 cursos ( 21,5 \% ) receberam o conceito 2 ; 48 cursos ( 8,4 \% ) o conceito 5 ; e 11 cursos ( 1,9 \% ) o conceito 1.
Nessa modalidade de oferta, 18 cursos ( 3,1 \% ) ficaram sem conceito ( SC ).
Dos 48 cursos participantes da modalidade a distância, o conceito 3 foi o valor modal, atribuído a 27 ( 56,3 \% ) cursos.
Dos demais cursos ofertados nessa modalidade, nove cursos ( 18,8 \% ) receberam conceito 4, nove cursos ( 18,8 \% ) receberam o conceito 2 ; um curso ( 2,1 \% ) recebeu o conceito 1 ; e nenhum curso ( 0,0 \% ) recebeu conceito 5.
Nessa modalidade de oferta, dois cursos ( 4,2 \% ) ficaram sem conceito ( SC ).
Na análise por região, observa -se, que, na região Norte, as instituições públicas participaram com quatro cursos ( 11,8 \% total regional ), dos quais nenhum ficou sem conceito ( SC ) ou recebeu os conceitos conceitos 5 e 1.
O conceito modal foi 4, atribuído a dois cursos, correspondendo a 50,0 \% dessa categoria na região.
O conceito 3 teve um curso ( 25,0 \% ), e o conceito 2 também teve um curso ( 25,0 \% ).
As instituições privadas participaram com 30 cursos na região Norte ( 88,2 \% do total regional ), e um curso ficou sem conceito ( SC ).
O conceito modal nas instituições privadas na região foi o 2, atribuído a 15 cursos, o que corresponde a 50,0 \% dessa categoria na região.
Os conceitos 3, 4 e 1 foram atribuídos a oito, quatro e dois cursos, respectivamente.
O conceito 5 foi atribuído a nenhum curso ( 0,0 \% ).
Nenhum curso ficou sem conceito ( SC ).
Dos cursos oferecidos na região Norte, 31 eram presenciais e um ficou sem conceito ( SC ).
Dos 31 cursos oferecidos, 16 obtiveram o conceito 2 ( modal ) ; oito o conceito 3 ; cinco o conceito 4 ; um o conceito 1 ; e nenhum o conceito 5.
Três cursos eram na modalidade a distância, tendo um recebido o conceito modal 4 ; um o conceito 3 ; um o conceito 1 ; nenhum o conceito 5 ; e nenhum o conceito 2.
Nenhum curso ficou sem conceito ( SC ).
Na região Nordeste, a rede privada concentra 132 dos 152 cursos participantes, o equivalente a 86,8 \% do total da região.
Entre os cursos oferecidos por IES privadas no Nordeste, 57 cursos foram avaliados com conceito 3, o conceito modal. 47 cursos obtiveram o conceito 2 ; 15 cursos o conceito 4 ; três cursos o conceito 1 ; e dois cursos obtiveram o conceito 5.
Nessa categoria administrativa e grande região, oito cursos ficaram sem conceito ( SC ).
As instituições públicas da região Nordeste participaram com 20 cursos ( 13,2 \% do total da região ).
Desses, 11 ( 55,0 \% ) obtiveram conceito 5, o valor modal, seguidos por nove cursos que receberam o conceito 4.
Os conceitos 3, 2 1 foram atribuídos a nenhum curso.
Nenhum curso ficou sem conceito ( SC ).
Dos cursos oferecidos na região Nordeste, 146 eram presenciais : sete ficaram sem conceito ( SC ), três obtiveram o conceito 1 ; 46 o conceito 2 ; 53 o conceito 3, o valor modal ; 24 o conceito 4 ; e 13 o conceito 5.
Seis cursos eram na modalidade a distância, tendo quatro cursos recebido o conceito modal 3 ; um o conceito 1 ; e nenhum curso recebeu os conceitos 4, 5 e 2.
Um curso ficou sem conceito ( SC ).
Na região Sudeste, a proporção de cursos da rede privada, 89,3 \%, correspondeu a 250 dos 280 cursos participantes.
Nessa combinação de categoria administrativa e grande região, o conceito modal foi o 3, atribuído a 119 cursos ; e seis cursos ficaram sem conceito ( SC ).
Os demais foram avaliados com os seguintes conceitos : 4 ( 68 cursos ), 2 ( 45 cursos ), 5 ( oito cursos ) e 1 ( quatro cursos ).
Dos 30 cursos oferecidos por instituições públicas na região Sudeste, 14 obtiveram o conceito 5 ; dez o conceito 4 ; quatro o conceito 3 ; dois obtiveram o conceito 2 e nenhum obteve o conceito 1.
Nenhum curso ficou sem conceito ( SC ).
Ainda na região Sudeste, grande parcela dos cursos ( 92,5 \% ) foi oferecida na modalidade presencial : 259 dos 280 cursos.
Para esses, a moda foi o conceito 3, com 110 cursos.
Os demais cursos presenciais da região foram \textbf{alocados} nos conceitos 4 ( 75 cursos ), 2 ( 42 cursos ), 5 ( 22 cursos ) e 1 ( quatro cursos ).
Seis cursos presenciais ficaram sem conceito ( SC ).
Na região Sudeste, 21 cursos foram oferecidos na modalidade a distância, não tendo cursos com conceito 1 e 5 ; cinco com conceito 2 ; 13 com conceito 3 ; e três com conceito 4 ;.
Nenhum curso ficou sem conceito ( SC ).
As instituições privadas concentraram 88 dos 100 cursos participantes da região Sul, 88,0 \% do total regional.
Desses, 42 cursos obtiveram conceito 3, o conceito modal.
Nessa combinação de categoria administrativa e grande região, 32 receberam o conceito 4 ; seis receberam o conceito 2 ; quatro o conceito 5 ; e um obteve o conceito 1.
Três cursos ficaram sem conceito ( SC ).
As instituições públicas na região Sul participaram com 12 cursos ( 12,0 \% dos cursos da região ), aos quais não houve cursos com os conceitos 1 e 2.
Ao conceito 3 foram atribuídos dois cursos, ao 4 cinco cursos e ao 5 cinco cursos.
Nenhum curso da região Sul oferecido por IES públicas ficou sem conceito ( SC ).
Na região Sul, 83 dos cursos foram oferecidos na modalidade presencial.
Para esses, a moda foi o conceito 3, com 36 cursos.
Os demais cursos presenciais da região foram \textbf{alocados} nos conceitos 4 ( 32 cursos ), 5 ( nove cursos ), 2 ( três cursos ) e 1 ( um curso ).
Ficaram sem conceito ( SC ) dois cursos.
Na modalidade a distância, foram oferecidos 17 cursos, sendo nenhum curso com conceito 1, três cursos com conceito 2, oito cursos com conceito 3 ( modal ), cinco cursos com conceito 4 e nenhum curso com conceito 5.
Um curso ficou sem conceito ( SC ) nessa região.
Na região Centro-Oeste, 49 dos 55 cursos participantes eram de IES privadas ( 89,1 \% em termos regionais ).
Desses, 19 cursos receberam conceito 3, o valor modal.
Dos demais cursos, 16 receberam o conceito 2, nove receberam o conceito 4, dois receberam o conceito 1 e um recebeu o conceito 5. dois ficaram sem conceito ( SC ).
Quanto aos seis cursos oferecidos por instituições públicas na região Centro-Oeste ( 10,9 \% do total regional ), nenhum ficou com conceitos 1 e 2 ; um ficou com conceito 3 ; dois com conceito 4 ; e três com conceito 5.
Nenhum ficou sem conceito ( SC ) nessa região.
Também na região Centro-Oeste, a maioria dos cursos foi ofertada na modalidade presencial : 54 ( 98,2 \% ).
Para esses, a moda foi o conceito 3, com 19 cursos ( 35,2 \% ).
Os demais cursos presenciais da região foram \textbf{alocados} nos conceitos 2 ( 16 cursos, 29,6 \% ), 4 ( 11 cursos, 20,4 \% ), 5 ( quatro cursos, 7,4 \% ) e 1 ( dois cursos, 3,7 \% ).
Dois cursos presenciais da região ficaram sem conceito ( SC ).
Na modalidade a distância, foi oferecido um curso, em que nenhum curso obteve conceito 1, 2, 4 ou 5.
Nenhum curso ficou sem conceito ( SC ) nessa região.
Para o único curso a distância nessa região, o conceito atribuído foi 3 ( modal ).
A título de síntese, no Gráfico 4.2, apresenta -se a distribuição do Conceito Enade, segundo a categoria administrativa da IES.
Nesse gráfico, como dito anteriormente, poligonais mais à esquerda representam Conceito Enade mais baixos, e poligonais mais à direita Conceito Enade mais alto.
É importante observar que, para cada poligonal presente nesse gráfico, a porcentagem constante no eixo y refere -se ao universo de cursos oferecidos por instituições públicas ( poligonal azul ) ou privadas ( poligonal verde ), respectivamente.
Os cursos em IES públicas ( poligonal azul ) apresentam uma poligonal mais à direita do que os em IES privadas ( poligonal verde ) e, consequentemente, uma distribuição de Conceito Enade com valores maiores. 4.3 CONCEITOS POR ORGANIZAÇÃO ACADÊMICA E GRANDE REGIÃO Na Tabela 4.3, encontra -se a distribuição dos conceitos atribuídos aos cursos participantes do Enade 2023, na área de Nutrição, por organização acadêmica, segundo as grandes regiões brasileiras.
Como mostram os dados da Tabela 4.3, dos 621 cursos de Nutrição participantes, 218 eram oferecidos em universidades, 262 em centros universitários e 138 em faculdades ; e 3 cursos foram oferecidos por CEFET/IF.
Essa distribuição corresponde a, respectivamente, 35,1 \%, 42,2 \%, 22,2 \% e 0,5 \% dos cursos.
De acordo com os dados apresentados na mesma tabela, de todos os 253 cursos avaliados com o conceito 3, 69 eram vinculados a universidades.
Os cursos vinculados a esse tipo de organização acadêmica tiveram o conceito 4 como conceito modal, atribuído a 77 cursos ( 35,3 \% ).
Os demais cursos ofertados por universidades receberam os conceitos 5 ( 40 cursos ), 2 ( 30 cursos ) e 1 ( um curso ).
Um curso ficou sem conceito ( SC ).
Dos cursos vinculados a centros universitários, o conceito modal foi o 3, atribuído a 132 cursos ( 50,4 \% ).
Nove cursos vinculados a esse tipo de organização acadêmica ficaram sem conceito ( SC ), e os outros cursos receberam os conceitos 4 ( 61 cursos ), 2 ( 49 cursos ), 1 ( sete cursos ) e 5 ( quatro cursos ).
Dos 138 cursos mantidos por faculdades, 53 ( 38,4 \% ) receberam conceito 2, o conceito modal.
O segundo conceito mais frequentemente atribuído aos cursos oferecidos por essa categoria de IES foi o conceito 3, recebido por 51 cursos.
Na sequência, 16 cursos obtiveram conceito 4 ; quatro cursos o conceito 5 ; e quatro cursos, o conceito 1.
Dos cursos oferecidos em faculdades, dez cursos ficaram sem conceito ( SC ).
Dos cursos oferecidos por CEFET/IF, o conceito modal foi o 4, atribuído a dois cursos ( 66,7 \% ).
Nenhum curso vinculado a esse tipo de organização acadêmica ficou sem conceito ( SC ), e um curso recebeu o conceito 3.
Nenhum curso foi atribuído aos conceitos 5, 2 e 1.
Considerando -se, separadamente, as regiões brasileiras, verifica -se que, na região Norte, as universidades participaram com oito cursos avaliados, aos quais foram atribuídos os seguintes conceitos : 1 a nenhum curso ; 2 a dois cursos ; 3, o conceito modal, a quatro cursos ; 4 a dois cursos ; e 5 a nenhum curso.
Nenhum curso ficou sem conceito ( SC ).
Os centros universitários da região Norte foram representados por 16 cursos, dos quais um curso obteve conceito 1 ; seis cursos o conceito 2 ; quatro cursos o conceito 3, o valor modal ; quatro cursos o conceito 4 ; e nenhum curso o conceito 5.
Um curso ficou sem conceito ( SC ).
As faculdades participaram com dez cursos na região Norte, os quais receberam os conceitos 1 ( um curso ), 2 ( oito cursos ), 3 ( um curso ), 4 ( nenhum curso ) e 5 ( nenhum curso ).
Nenhum curso ficou sem conceito ( SC ).
Nos CEFET/IF, por não haver nenhum curso, não há classificação dos conceitos para essa grande região.
Na região Nordeste, as universidades participaram com 31 dos 152 cursos da área de Nutrição oferecidos.
O conceito modal foi 5, atribuído a 11 cursos.
Outros 11 cursos obtiveram o conceito 4 ; oito cursos o conceito 3 ; um curso o conceito 2 ; e nenhum curso o conceito 1.
Nenhum curso ficou sem conceito ( SC ). cursos o conceito 3 ; um curso o conceito 2 ; e nenhum curso o conceito 1.
Nenhum curso ficou sem conceito ( SC ).
Os centros universitários contaram com 69 cursos participantes na região Nordeste, 35 dos quais receberam o conceito modal 3.
Para os demais cursos, foram atribuídos os conceitos 2 ( 18 cursos ), 4 ( oito cursos ), 1 ( dois cursos ) e 5 ( um curso ).
Cinco cursos oferecidos por centros universitários no Nordeste ficaram sem conceito ( SC ).
As faculdades foram representadas por 51 cursos na região Nordeste, os quais receberam os conceitos 1 ( um curso ), 2 ( 28 cursos ), 3 ( 14 cursos ), 4 ( quatro cursos ) e 5 ( um curso ).
Três cursos ficaram sem conceito ( SC ).
Na região Nordeste, os CEFET/IF participaram com um curso, que foi atribuído ao conceito 4.
Na região Sudeste, as universidades concentraram 121 dos 280 cursos de Nutrição da região.
Dos cursos oferecidos em IES com este tipo de organização acadêmica, nessa região, o conceito modal foi 3, atribuído a 41 cursos.
Os demais cursos receberam os conceitos 4 ( 37 cursos ), 2 ( 23 cursos ), 5 ( 18 cursos ) e 1 ( um curso ).
Um curso ficou sem conceito ( SC ).
Os centros universitários participaram com 108 cursos na região Sudeste, dos quais 59 obtiveram o conceito modal 3.
O conceito 4 foi atribuído a 30 cursos ; o conceito 2 a 14 cursos ; o conceito 1 a três cursos ; e o conceito 5 a um curso.
Um curso ficou sem conceito ( SC ) nessa combinação de organização acadêmica e região.
Já as faculdades foram representadas por 50 cursos na região Sudeste, dos quais quatro ficaram sem conceito ( SC ).
O conceito modal 3 foi atribuído a 23 cursos.
Aos demais cursos foram atribuídos os conceitos 4 ( dez cursos ), 2 ( dez cursos ), 5 ( três cursos ) e 1 ( nenhum curso ).
Os CEFET/IF participaram com um curso na região Sudeste, que foi atribuído ao conceito 4.
Dos 100 cursos da região Sul, 44 eram mantidos por universidades, para os quais o conceito modal foi 4, atribuído a 21 cursos.
Os demais cursos receberam os conceitos 3 ( 12 cursos ), 5 ( oito cursos ), 2 ( três cursos ) e 1 ( nenhum curso ).
Em IES com esse tipo de organização acadêmica, nenhum curso da região Sul ficou sem conceito ( SC ).
Dos 43 cursos participantes de centros universitários da região Sul, dois ficaram sem conceito ( SC ), e 23 cursos receberam conceito 3, o valor modal.
Os conceitos restantes ficaram assim distribuídos : 4 ( 14 cursos ), 2 ( três cursos ), 5 ( um curso ) e 1 ( nenhum curso ).
Quanto aos 13 cursos vinculados a faculdades na região Sul, nove cursos receberam o conceito 3, valor modal.
Os demais conceitos ficaram assim distribuídos : 4 ( dois cursos ), 1 ( um curso ), 5 ( nenhum curso ) e 2 ( nenhum curso ).
Dos 13 cursos, um ficou sem conceito ( SC ).
Na região Centro-Oeste, 14 dos 55 cursos oferecidos eram mantidos por universidades.
Para esse tipo de organização acadêmica, o conceito modal foi 4, atribuído a seis cursos.
Em relação às universidades, nenhum curso da região ficou sem conceito ( SC ).
Os demais cursos receberam os conceitos 3 ( quatro cursos ), 5 ( três cursos ), 2 ( um curso ) e 1 ( nenhum curso ).
Os centros universitários da região Centro-Oeste contaram com 26 cursos : 11 cursos foram avaliados com o conceito 3 ( modal ) ; oito cursos com o conceito 2 ; cinco cursos com o conceito 4 ; um curso com o conceito 5 ; e um curso obteve o conceito 1.
Nenhum curso ficou sem conceito ( SC ).
Dos 14 cursos oferecidos em faculdades na região Centro-Oeste, sete cursos receberam o conceito 2 ( modal ).
Os conceitos restantes ficaram assim distribuídos : 3 ( quatro cursos ), 1 ( um curso ).
Não obtiveram cursos os conceitos 5 e 4.
Nessa região, dois cursos oferecidos em faculdades ficaram sem conceito ( SC ).
No que se refere aos CEFET/IF da região Centro-Oeste, um curso foi avaliado e atribuído ao conceito 3.
De maneira semelhante à forma como se procedeu em relação à seção anterior, esta seção encerra -se com o Gráfico 4.3, que apresenta a distribuição do Conceito Enade dos cursos de Nutrição segundo a organização acadêmica da IES.
No Gráfico 4.3, apresenta -se a distribuição do Conceito Enade segundo a organização acadêmica da IES.
A leitura desse último gráfico é semelhante à dos anteriores, seguindo -se a lógica de que poligonais deslocadas para a esquerda correspondem aos conceitos mais baixos ( 1, 2 e 3 ) e poligonais deslocadas para a direita correspondem a Conceitos Enade mais altos.
Nas universidades ( linha azul ), todos os cursos obtiveram conceito 5 ou abaixo disso, já que a poligonal cumulativa alcança 100 \% para o conceito 5.
Além disso, os cursos de universidades apresentam uma poligonal mais à direita do que as demais, denotando uma melhor distribuição de Conceitos Enade.
No outro extremo, encontram -se os cursos mantidos por faculdades ( linha vermelha ), apresentando as respectivas poligonais mais à esquerda e, consequentemente, uma distribuição de Conceitos Enade com os menores valores.
Os cursos mantidos por centros universitários ( linha verde ) e CEFET/IF ( linha laranja ) apresentam poligonais mais centrais e, consequentemente, uma distribuição de Conceitos Enade com os valores intermediários.

% matched lemmas: alocar
\end{textsample}
